\documentclass[11pt,a4paper, final]{moderncv}

\usepackage{color}
\usepackage{fontspec}
\moderncvtheme{classic}

% the (custom) color which will be used in the cv
\definecolor{color1}{RGB}{1, 52, 64}

% scale the page layout
\usepackage[scale=0.95]{geometry}

% change width of the column with the dates
\setlength{\hintscolumnwidth}{2.5cm} 

% required when changing page layout lengths
\AtBeginDocument{\recomputelengths} 
\usepackage{xunicode}
\usepackage{xltxtra}
\usepackage[utf8]{inputenc}

% german word break/hyphenation rules
\usepackage[ngerman]{babel}

% insert dummy text (used in the letter)
\usepackage{lipsum} 

% used for \begin{comment}...\end{comment}
\usepackage{verbatim} 

% use guilllemets in bibliography
\usepackage[babel,german=guillemets]{csquotes}

\usepackage[
	sorting=none,
	minbibnames=8,
	maxbibnames=9,
	block=space
]{biblatex}

\bibliography{publications}


% get rid of the number-labels ([1], [2], etc.) in front of publications
\defbibenvironment{midbib}
{\list
	{}
	{
		\setlength{\leftmargin}{0mm}
		\setlength{\itemindent}{-\leftmargin}
		\setlength{\itemsep}{\bibitemsep}
		\setlength{\parsep}{\bibparsep}}
	}
	{\endlist}
{\item}


% add [DOI] and [PDF] fields at the end of each publication entry
\DeclareSourcemap{
	\maps[datatype=bibtex]{
		% the bibtex entry 'mydoi' gets mapped to 'usera'
		\map{
			\step[fieldsource=mydoi]
			\step[fieldset=usera, origfieldval]
		}
		% the bibtex entry 'mypdf' gets mapped to 'usera'
		\map{
			\step[fieldsource=mypdf]
			\step[fieldset=userb, origfieldval]
		}
	}
}

% [DOI] entries in publication
\DeclareFieldFormat{usera}{\color{color1}[\href{#1}{\textsc{doi}}]}
\AtEveryBibitem{
	% put [DOI] stuff at the end of a publication entry
	\csappto{blx@bbx@\thefield{entrytype}}{% 
		\iffieldundef{usera}{ 
			% this gets invoked, once nothing is supplied
			% via the mypdf or mydoi value.
			% you could e.g. display a default thing here.
		}{\space\printfield{usera}}}
}

% [PDF] entries in publication
\DeclareFieldFormat{userb}{\color{color1}[\href{#1}{\textsc{pdf}}]}

\AtEveryBibitem{
	% put [DOI] stuff at the end of a publication entry
	\csappto{blx@bbx@\thefield{entrytype}}{\iffieldundef{userb}{}{\printfield{userb}}}
}

\renewcommand*{\mkbibnamegiven}[1]{%
	\ifitemannotation{highlight}
	{\textbf{#1}}
	{#1}}

	\renewcommand*{\mkbibnamefamily}[1]{%
	\ifitemannotation{highlight}
	{\textbf{#1}}
	{#1}
}


% Minion Pro is used as the main font, if you don't
% have it installed uncomment this line or choose another pretty font
\setmainfont[Numbers=OldStyle]{Times New Roman}

% the moderncv package will populate a lot of the pdf meta-information.
% this can be used to change some of these infos.
\AfterPreamble{\hypersetup{
  pdfcreator={XeLaTeX},
  pdftitle={Fictional CV of Hai Jiang}
}}

% for the icons (telephone, globe). i found the icons provided by the
% fontawesome package prettier than the standard moderncv icons.
\usepackage{fontawesome5}

% personal data
\firstname{Hai}
\familyname{Jiang}
\address{Baoan District, Shenzhen, China}{}

% i use the extrainfo command for additional information, since i 
% want to use custom icons and have finer control over spacing.
\extrainfo{
	\faPhone\hspace{0.3em}+86 19867713757\\
	{\small\faEnvelope}\hspace{0.3em}pigejianghai97@gmail.com%\\
	% \faGlobe\hspace{0.3em}http://isaac.asimov.com
}

% picture, resized to a height of 84pt
\photo[134pt]{picture/XFOB5134}

% spacing before (sub)sections
\newcommand{\spacesection}{\vspace{0.4cm}}
\newcommand{\spacesubsection}{\vspace{0.2cm}}


%===========================
\begin{document}

\recipient{Motivation Letter}{}%{Shenzhen\\Guangdong\\China} 
\date{\today}
\opening{To whom it may concern,} %Dear Prof. Dr. Hussam Amrouch,
\closing{Sincerely,}
\enclosure[Attached]{curriculum vit\ae{}}

\makelettertitle
I am writing to apply for the PhD posotion focusing on the Big Model in the Artificial Intelligence at your university.
\ \\
\ \\
With a deep-seated passion for Computer Science and Mathematics, 
I began my studies in Computational Mathematics in 2019 after earning my Bachelor's degree in Information Security in China.
While focusing on Deep Learning technology, I participated in a research project aimed at imitating personal handwriting styles.
The project involved learning the Chinese handwritten style of Shiing Shen Chern from approximately 220 characters.
Through this experience, I gained knowledge of Generative Adversarial Networks 
and developed skills in manuscript searching, data preprocessing, and code replication.
As a result, I completed my Master's thesis on the personalized imitation of Chinese handwritten characters using GANs.
\ \\
\ \\
Furthermore, to integrate Deep Learning with numerical approximation theory, 
I participated in the Compressed Sensing MRI ADMM-Net project.
Traditional numerical algorithms perform well on simple images, 
but they fail to preserve detailed information after numerous iterations.
With advancements in Deep Learning, 
many researchers have employed Convolutional Neural Networks to unfold equations 
and automatically adjust parameters during iterations to preserve details. 
Through this project, I gained knowledge of convergent algorithm theory and proofs, 
as well as the integration process with Deep Learning. 
This experience significantly enhanced my skills in both programming and theoretical understanding.
\ \\
\ \\
After earning my Master's Degree in Computational Mathematics, 
I worked as a Research Assistant in a group at Sun Yat-sen University dedicated to Computer-Assisted Intervention. 
During my tenure, I participated in proposal writing 
for a key program from the Science and Technology of China (2023YFE0204300, which was successfully approved). 
I also completed two progress reports and the completion report for projects funded 
by the National Natural Science Foundation of China (Grants 81971691 and 12126610). 
Additionally, I wrote the technological specifications for a patent application and a medical device application. 
Over the past two years, I have contributed to three projects focused on 
diagnosing Placenta Accreta Spectrum Disorders, 
predicting metastasis in Axillary Sentinel Lymph Nodes, 
and predicting the response to Neoadjuvant Chemotherapy based on MRI. 
The first two projects yielded excellent experimental results but lacked technological novelty. 
Throughout these tasks, I consolidated skills in programming (Python, PyTorch, and TensorFlow) and manuscript reading.
\ \\
\ \\
I believe that the PhD position in your research unit will be a life-changing experience to 
achieve the career I aspire for. 
Your research field the Machine Learning Mechanism are aligned my core interests. 
I am genuinely committed to investigating pragmatic ways to 
maximize use of Engineering and Computer Science. 
% On the contrary, the 
\ \\
\ \\
I am confident that I can thrive as a promising candidate for this PhD position. 
My academic background and program experience, along with my personal skills, 
demonstrate my compatibility with this program. 
In summary, I have acquired extensive knowledge and skills in Computer Science and Computational Mathematics. 
Additionally, my research in Deep Learning, Medical Science, 
and Mathematics has equipped me to be a competent team player.
\ \\
\ \\
Until then, I am grateful for the opportunity to apply to this institution and I am 
ready to discuss my future with you should my candidacy prove successful.

\makeletterclosing

\clearpage
%%%%%%%%%%%%%%%%%%%%%%%%%%%%%%%%%%%%%%%%%%%%%%%%

\maketitle
\section{Personal Information}
\cvline{Name}{Hai Jiang}
\cvline{Birth}{21.Dec.1997}
% \cvline{Geburtsort}{Petrovichi, Russia}
% \spacesection
\section{Educational Background}
\cventry{\textsc{2019--2022}}{M.Sci. Computational Mathematics}{Nankai University}{Tianjin}{\emph{China}}{}
\cvline{Related Courses}
{Matrix Computation, Appriximation theory and methods, Numerical Optimization, 
Functional Analysis, Convex Analysis, Variational Analysis, 
Real and Complex Analysis, Computational Geometry, Digital Image Processing, Fourier Analysis}
\spacesection
\cventry{\textsc{2014--2018}}{B.Eng. Information Security}{Lanzhou University}{Lanzhou}{\emph{China}}{}
\cvline{Related Courses}
{Advanced Mathematics, Linear Algebra, Elementary Number Theory, Statistic and Probability Theory, 
Information Theory, Discrete Mathematics, Operating System, Data Structure, 
C and C++ Programming Laboratory, Java Programming Laboratory, 
Database Theory and Laboratory, Computer Organization and Design, Circuit Analysis, }
% \cventry{\textsc{1935--1939}}{B.Sc. Chemistry}{University Extension}{New York}{\emph{China}}{}
% \cventry{\textsc{2011--2014}}{High School}{Songgang High School}{Shen Zhen}{\emph{China}}{}
% \spacesection
\section{Research and Project Experience}
\subsection{Space-squeeze method for multi-class classification of metastatic lymph nodes of breast cancer}
\cventry{\textsc{12.2023--01.2024}}
{Advisor: Yao Lu, Xiang Zhang, Cooperator: Gary Pan}
{Sun Yat-sen University}
{Submitted to MICCAI2024}{Early Reject}
{
■ A task under the National Natural Science Foundation of China (NSFC) project, Grant 81971691, has been completed.
This project, which focused on breast cancer, 
aimed to develop a prediction model for the Chinese female population using full-field digital mammography and ultrasonography.\\
■ The task, based on the Dual-Energy CT (DECT) image sequence in the coronal position with sentinel Axillary Lymph Node (ALN) annotation, 
aimed to classify the metastatic status of ALNs. 
We proposed two technologies: 
firstly, a feature space squeezing method to explore the correlation of energy levels; 
secondly, a representation space squeezing method by introducing virtual classes, 
thereby compacting the intra-class space and sharpening the inter-class boundaries.\\
■ I gained my first complete research experience through this task, 
which included conducting literature research, designing experiments, and writing papers.
}
\subsection{Artificial Intelligence Model For MRI Prenatal Classification of Placenta Accreta Spectrum Disorders}
\cventry{\textsc{11.2022--08.2023}}
{Advisor: Yao Lu, Ting Song, Cooperator: Qiongting Liu}
{Sun Yat-sen University}
{Submitted to American Journal of Roentgenology}{Rejected}
{
■ A task of the China Department of Science and Technology Key Grant 2023YFE0204300 is currently in progress. 
This project focuses on breast cancer within the clinical setting. Guided by the expertise of clinical radiologists, 
the project aims to develop models with clinical interpretability and generalization. 
It seeks to address four key challenges in clinical diagnosis: 
small sample learning, multi-modality, interpretability, and imbalanced data distribution.\\
■ Based on the T2-weighted MRI images in the sagittal position and the annotations of the placenta and serous line, 
the task aims to classify the subtype (PA: placenta accreta, PI: placenta increta, PP: placenta percreta) 
of patients with biopsy-proven PAS positivity. 
I proposed two shallow neural networks for prediction: 
one for the classification of PP/non-PP and another for the classification of PA/PI. 
The first model achieved an AUC (Area Under the ROC Curve) of 0.94, while the second achieved an AUC of 0.85.\\
■ Gained experience in literature research, data preprocessing, and model building..
}
% \subsection{Multi-view based Atypical Architectural Distortion Detection of Breast Cancer}
% \cventry{\textsc{08.2022--03.2023}}
% {Advisor: Yao Lu, Cooperator: Gary Pan}
% {Sun Yat-sen University}
% {Submitted to the Journal of Physics in Medicine and Biology}{Accepted}
% {
% ■	A task of the project National Natural Science Foundation of China(NSFC) under Grant 12126610, is in progress. 
% The project aimed to develop a personalized prediction model for early breast cancer prevention in the Chinese female population. \\
% ■	The task aims to detect the breast architectural distortion based on the DBT(Digital Breast Tomosynthesis) volume with two different angles. 
% Inspired by the Siamese-Networks initially used in face recognition, the task utilized the Siamese architecture to extract and compare information from both views. 
% Experimental results show that the corresponding triplet module did help the False-Positive reduction from two angles.\\
% ■	My involvement in the project focused on preparatory works including providing ideas and consultation.\\
% ■	I gained my knowledge in multiple machine learning methods, 
% including automatic assessment of architectural distortion, detection technology in deep learning, et al.\\
% }

\spacesubsection
\section{Experiments Experience During Master}
\subsection{Deep ADMM-Net for Compressed-Sensing MRI}
\cventry{\textsc{01.2022--06.2022}}{Advisor: Prof. Louis Wu}{Nankai University}
{}{}
{
■ ADMM-model from the manuscript “Deep ADMM-Net for Compressed-Sensing MRI.”\\
■ The goal of Compressed Sensing (CS) is to recover the original image from a subset of k-space (i.e., Fourier Space) data. 
Such ill-posed inverse problems can be addressed by various methods, 
including primal-dual methods and ADMM (Alternating Direction Method of Multipliers). 
While traditional methods perform well in theory and with simple images, 
they often fail to preserve detailed information after numerous iterations. 
With advancements in Neural Networks and Deep Learning, many researchers are leveraging AI to unfold the equations, 
adjusting numerical values during the iterations to better preserve image details.\\
■ Gained my second experience in numerical experiments, combining mathematical equations with modern deep learning technology. 
Subsequently, I completed the numerical experiments, 
reproducing the mathematical iteration equations in both C++ and Python, 
and implemented them in the PyTorch framework.
}
\subsection{GANs based Personal style imitation of Chinese handwritten characters}
\cventry{\textsc{09.2021--02.2022}}
{Advisor: Prof. Louis Wu, Yunhua Xue}
{Nankai University}
{}{}
{
■ The task was to learn the Chinese handwritten style of Shiing Shen Chern from approximately 220 pictures of his handwritten characters, 
presented on the occasion of his 110th birth anniversary. 
Utilizing GANs (Generative Adversarial Networks), I employed CycleGANs (Cycle Consistent Generative Adversarial Networks) 
and relevant Cycle-Consistent Loss for style-transfer learning. 
Subsequently, I proposed the Deformable Generative Network, based on the deformable convolutional operator, 
to further enhance the learning of the handwritten style.\\
■ I gained my first complete research experience, encompassing data collection, data pre-processing, 
network building, network training, and results comparison. \\
% As a result, the idea of two networks sounds good in theory, but works much differently in practice. 
% The CycleGANs work much better than the other one.\\
■ I completed my Master thesis, titled "Intimated Personalized Style of Chinese Handwritten Characters Based on GANs".
}
\subsection{Total-Variation Denoise ROF-model}
\cventry{\textsc{01.2021--04.2021}}{Advisor: Prof. Yunhua Xue}{Nankai University}
{}{}
{
■ ROF-model from the manuscript “Nonlinear Total Variation Based Noise Removal Algorithms.”\\
■ The manuscript proposes a new constraint, also known as a penalty term, namely Total Variation (TV), 
for recovering the original image from a denoised (salt and pepper noise added manually) one. 
The mathematical equations converge to a stable point after numerous iterations, 
which can be solved using the Euler-Lagrange Equations. 
When comparing their method to the Mean Square Error (MSE) Loss Function, 
the numerical results show that the TV-term performs better in preserving boundaries.\\
■ Gained my initial experience in numerical experiments, utilizing both C++ and Python, as described in the manuscript.
% As the classical denoise algorithm, the ROF model did well in protecting boundaries while lacking smoothness in the internal. 
% Modern deep-learning methods are much more powerful tools for maintaining smoothness remaining or boundary protection in image denoise.
}

% \newpage
\spacesection
\section{Personal Skills}
% \subsection{Personal Skills}
\cvline{Language Proficiency}
{Mandarin(Native), Cantonese; CET6 476/710; CET4 544/710; IELTS 6.5.}
\cvline{Programming}
{C/C++, Python 3, MATLAB, PyTorch, Tensorflow, Latex.}
\cvline{Self-study}
{A good self-study capability in science or engineering, including mathematics, computer science.}
% \spacesection
% \subsection{Computer Sciense, Information Security}
% \cvline{Related Courses}
% {Advanced Mathematics, Linear Algebra, Elementary Number Theory, 
% Statistic and Probability Theory, Information Theory, Discrete Mathematics, 
% Operating System, Data Structure, Standard C&C++ Programming Laboratory, 
% Standard Java Programming Laboratory, Database Theory and Laboratory, Computer Organization and Design}
% {Application Realm}{Digital Image Processing, MRI}
% \cvcomputer{Consectetuer}{Felis, Adipiscing, Vitae}
%            {Donec}{Felis, Erat, Congue non}
% \spacesection
% \subsection{Computer Science, Information Security}
% \cvcomputer{Related Courses}
% {Advanced Mathematics, Linear Algebra, Elementary Number Theory, 
% Statistic and Probability Theory, Information Theory, Discrete Mathematics, 
% Operating System, Data Structure, Standard C&C++ Programming Laboratory, 
% Standard Java Programming Laboratory, Database Theory and Laboratory, Computer Organization and Design}
% % {}{}
% {Application Realm}{Information Security}
% \cvcomputer{Eleifend}{Morbi blandit, Ligula}
%            {Accumsan}{Nam ipsum, suscipit}

\spacesection
\section{References}
\subsection{Prof. Yao Lu}
\cvline{}
{Research assistant at the Computational Medical Image Laboratory at Sun Yat-sen University.}
\cvline{Name}{Yao Lu}
\cvline{Work Address}
{Xingzhen Building, Sun Yat-sen University, Panyu District, Guangzhou, Guangdong Province, China, Room B401}
\cvline{Tel}{+86 18664509325}
\cvline{E-mail}{luyao23@mail.sysu.edu.cn}
\subsection{Prof. Qing Ling}
\cvline{}
{Research assistant at the Computational Medical Image Laboratory at Sun Yat-sen University.}
\cvline{Name}{Qing Ling}
\cvline{Work Address}
{Xingzhen Building, Sun Yat-sen University, Panyu District, Guangzhou, Guangdong Province, China, Room B408}
\cvline{Tel}{+86 18654109710}
\cvline{E-mail}{lingqing556@mail.sysu.edu.cn}
% \cvline{Prof. Yao Lu}{
% 	Lorem ipsum dolor sit amet, consectetuer adipiscing elit. Ut purus
% 	elit, vestibulum ut, placerat ac, adipiscing vitae, felis. 
% }
% \cvline{Vestibulum}{
% 	Maecenas lacinia. Nam ipsum ligula, eleifend at, accumsan nec, suscipit a,
% 	ipsum. Morbi blandit ligula.
% }


% \spacesection
% \section{Miscellaneous}
% \cventry{\textsc{1958--1963}}{Adipiscing Elit}{New York}{}{}{
% 	Vivamus viverra fermentum felis. Donec nonummy pellentesque ante. Phasellus
% 	adipiscing semper elit. Pellentesque cursus luctus mauris.
% }


\spacesection
\section{Academic Writings}
% {\color{color1}\textsc{English}}
% \ \\
1. Space-squeeze method for multi-class classification of metastatic lymph nodes of breast cancer, MICCAI2024, Early Reject
\ \\
2. GANs based Personal style imitation of Chinese handwritten characters, Master thesis(in Chinese)

% \nocite{*}
% \printbibliography[heading=none, env=midbib, keyword=academic]

% \ \\

% {\color{color1}\textsc{Fictional Writings}}
% \printbibliography[heading=none, env=midbib, keyword=fiction]



%=============================
% this part is a simple cover letter
\clearpage
\end{document}
